\documentclass[11pt]{letter}
\usepackage[margin=2.5cm]{geometry}
\usepackage{hyperref}

\signature{Joel Ayala-Baez}
\address{Independent Researcher\\jayalabaez@gmail.com}

\begin{document}
\begin{letter}{The Editor\\Monthly Notices of the Royal Astronomical Society}

\opening{Dear Editor,}

I am submitting the manuscript entitled \textbf{``HD~338492: A
High-Mass Dark Companion Candidate from Gaia DR3
Astrometric--Spectroscopic Binary Data''} for consideration for
publication in MNRAS.

This paper presents an archival multi-wavelength analysis of
HD~338492, a single-lined spectroscopic binary in the Gaia~DR3
non-single-star catalogue whose dynamical properties point to a
dormant black hole (BH) companion.  The key results are:

\begin{itemize}
\item A spectroscopic mass function $f(M) = 2.39\,M_\odot$, which
  exceeds the Chandrasekhar limit from dynamics alone.
\item A minimum companion mass $M_{2,\min} = 6.62\,M_\odot$ (for
  $M_1 = 4.40\,M_\odot$), well above the neutron-star ceiling.
\item No luminous companion detected in any photometric band
  (optical through mid-infrared), and no X-ray emission.
\item Four of six alternative non-BH scenarios excluded on dynamical
  and photometric grounds; the remaining two (hierarchical triple
  and stripped helium star) are strongly constrained but require
  targeted follow-up to close definitively.
\end{itemize}

\textbf{What this paper does, and what it does not claim.}
This work identifies HD~338492 as a strong BH \emph{candidate}
based solely on publicly available archival data.  We are
transparent about its limitations: the stellar characterisation
relies on the SIMBAD B9 spectral classification (not independently
verified), the SED uses a blackbody approximation, and no
independent radial-velocity confirmation exists.  We present a
sensitivity analysis (Table~4) demonstrating how the BH probability
depends on the adopted mass threshold, primary-mass prior, and
extinction.  Under fiducial assumptions $P(\mathrm{BH}) = 99.7\%$,
but this drops to $\sim\!89\%$ if the B9 anchor is abandoned
entirely --- underscoring the importance of follow-up
spectroscopy.  We explicitly identify the follow-up programme
needed to close the remaining gaps: independent multi-epoch
radial velocities, a modern spectroscopic classification, and UV
observations.

The discovery of Gaia~BH1, BH2, and BH3 has demonstrated the
power of Gaia astrometric--spectroscopic solutions for finding
dormant BHs.  These landmark systems were each confirmed with
independent RV campaigns.  By presenting a carefully vetted
archival candidate with a clearly defined follow-up programme, this
paper aims to expand the search space for dormant BHs from DR3 and
to motivate the observational resources needed for confirmation.

The full analysis pipeline (six reproducible Python scripts),
input data, and all intermediate results are publicly available at
\url{https://github.com/jayalabaez/hd338492-black-hole-candidate}
and will be archived on Zenodo at proof stage.

No portion of this work has been published or is under review
elsewhere.  The manuscript is approximately 5\,000 words in
length with 6 figures and 4 tables.

I suggest the following potential referees with expertise in Gaia
compact-object binaries:
\begin{itemize}
\item Kareem El-Badry (Caltech) --- lead author of Gaia BH1--BH3
\item Tomer Shenar (University of Amsterdam) --- binary-star
  spectroscopy and BH candidates
\item Tsevi Mazeh (Tel Aviv University) --- Gaia spectroscopic
  binaries
\end{itemize}

Thank you for considering this submission.

\closing{Sincerely,}
\end{letter}
\end{document}
